\documentclass[11pt, a4paper]{ctexart}

\usepackage{geometry}
\geometry{a4paper, margin=1in}
\usepackage{graphicx}
\usepackage{booktabs}
\usepackage{float}
\usepackage{hyperref}
\usepackage{amsmath}
\usepackage{caption}

\title{\textbf{基于结构与语义特征的数学证明来源判别研究}\\ \Large 人类解答与多种大语言模型解答的分类器构建}
\author{wjj}
\date{\today}

\begin{document}

\maketitle

\begin{abstract}
本报告围绕“数学证明文本来源判别”这一任务，系统介绍了一个用于区分人类解答与四种大语言模型解答（Deepseek、Kimi、GLM、Qwen）的多分类器构建过程。与直接使用黑盒深度模型不同，本项目采用“可解释特征工程 + 传统机器学习分类器”的路线：一方面使用 TF-IDF 捕捉词汇层面的局部模式，另一方面设计 28 个与数学证明写作强相关的结构、公式与语义特征，以显式建模不同来源在段落组织、LaTeX 使用方式、推理连接词和祈使句习惯上的差异。在固定 200 道题测试集的条件下，我们进一步考察了训练题量从 20 道到 800 道的变化对五种分类器性能的影响。实验表明，树集成模型在这一任务上表现显著优于线性模型，其中 Hist Gradient Boosting 在 800 道训练题设置下达到 95.5\% 的测试准确率，说明数学证明来源中确实存在可稳定提取的“生成指纹”。
\end{abstract}

\section{引言}
随着大语言模型在数学推导、证明书写和 LaTeX 表达方面的能力快速提升，单纯从肉眼阅读角度判断一份数学证明是由人类独立完成还是由模型生成，已经变得越来越困难。这一问题不仅与学术诚信直接相关，也为“风格溯源”“机器生成文本识别”“跨模型写作差异研究”等方向提供了一个非常具体且有挑战性的应用场景。

本项目的目标并不是判断一段文本“像不像 AI”，而是更进一步：在同一道数学题上，识别一份解答究竟来自五类来源中的哪一类，即 Human、Deepseek、Kimi、GLM 或 Qwen。为达成这一目标，我们希望构建的是一个同时满足以下三个要求的分类系统：
\begin{itemize}
    \item \textbf{准确}：在多分类场景下保持较高识别率；
    \item \textbf{可解释}：能够说明模型究竟依据哪些结构性证据作出判断；
    \item \textbf{可扩展}：当训练样本规模变化时，依然能稳定分析不同算法的表现。
\end{itemize}

因此，本项目没有直接采用大型预训练编码器作为最终方案，而是选择了一条更利于分析的技术路线：通过构造双轨特征流水线，将“词汇信息”和“数学证明结构信息”显式编码，再交由多种经典机器学习模型进行训练与比较。

\section{方法（Methods）}
本节从“数据收集与处理”“特征提取”“模型训练”三个层面详细说明分类器是如何从头构建出来的。

\subsection{数据收集与处理}

\subsubsection{原始数据来源}
项目的数据基础来自 HuggingFace 上的 \texttt{StackMathQA} 数据集。我们从中选取了 1000 道纯英文数学题目，每道题目都自带一个高质量的人类解答。此时，每一道题至少对应两部分核心信息：题干（problem）和人类标准解答（human）。

\subsubsection{多模型答案生成}
为了将该任务转化为五分类问题，我们对同一批题目分别调用四个大语言模型接口生成对应解答：
\begin{itemize}
    \item Deepseek-Chat
    \item Kimi / Moonshot
    \item GLM-4
    \item Qwen-Plus
\end{itemize}

生成时统一使用英文提示词，并明确要求输出严格、完整、偏数学证明风格的英文答案，尽量采用标准 LaTeX 公式书写。这样做的目的是尽可能将“内容主题差异”控制住，让分类器主要学习“表达风格差异”而不是“题目差异”。

\subsubsection{数据组织方式}
最终，每一道题目都对应五个来源版本：
\[
\text{problem} \rightarrow \{\text{human}, \text{deepseek}, \text{kimi}, \text{glm}, \text{qwen}\}
\]

因此总数据量为：
\[
1000 \text{ questions} \times 5 \text{ sources} = 5000 \text{ records}
\]

每条记录包含唯一题目编号、题干文本、答案文本以及来源标签。这样的设计有两个重要优点：
\begin{enumerate}
    \item 同一道题在五种来源之间形成天然配对，便于比较同题不同解法的写作风格；
    \item 训练集/测试集可以按题号切分，而不是按答案随机切分，从根本上避免同一道题的不同来源版本同时出现在训练与测试中造成数据泄漏。
\end{enumerate}

\subsubsection{训练集与测试集划分原则}
在后续实验中，我们始终先按题目编号进行划分，再将对应的五条答案整体放入训练集或测试集。这样保证：
\begin{itemize}
    \item 测试集中的题目在训练阶段完全不可见；
    \item 分类器不能依靠“记住某道题的内容”来预测来源；
    \item 最终结果更能反映模型对写作风格、结构特征和公式习惯的真实泛化能力。
\end{itemize}

\subsection{特征提取}

\subsubsection{总体思路：双轨特征流水线}
数学证明文本与普通自然语言文本不同，其来源差异不仅存在于词汇选择上，更存在于段落组织、公式嵌入方式、逻辑推进节奏和 LaTeX 使用习惯上。因此我们设计了一个双轨特征提取框架，并在 \texttt{sklearn} 中通过 \texttt{ColumnTransformer} 将其统一组合：
\begin{enumerate}
    \item \textbf{TF-IDF 词袋轨道}：提取局部词汇和短语模式；
    \item \textbf{自定义结构特征轨道}：提取数学证明写作中的全局结构与风格指纹。
\end{enumerate}

\subsubsection{TF-IDF 轨道}
第一条轨道使用 \texttt{TfidfVectorizer} 对答案文本进行向量化，设置上主要包括：
\begin{itemize}
    \item 最大特征维度约 800--1000；
    \item 使用 unigram + bigram；
    \item 保留普通单词与 LaTeX 宏符号；
    \item 不依赖主题语义建模，而强调局部词汇分布差异。
\end{itemize}

该轨道主要负责捕捉如下模式：
\begin{itemize}
    \item 某些模型偏好的固定衔接短语；
    \item 习惯性的结论表达方式；
    \item 常见 LaTeX 宏和证明起手式的使用频率。
\end{itemize}

但是，TF-IDF 只能看到“有哪些词”，看不到“这些词是如何被组织起来的”。因此它不能单独承担整个分类任务。

\subsubsection{自定义结构特征轨道}
第二条轨道由我们手工构造的 \texttt{TextFeatureExtractor} 完成。它通过正则表达式和基础统计，对每篇答案计算 28 个可解释特征。之所以称其为“物理特征”，是因为它们直接描述了一份数学证明在页面上、在段落中、在公式组织中的可观察行为，而不依赖隐式语义表示。

这些特征可以分为三类：

\paragraph{1. 宏观结构特征}
这一组特征试图回答“作者是如何组织整篇证明的”：
\begin{itemize}
    \item 文本总长度（length）
    \item 总行数（num\_lines）
    \item 平均行长（avg\_line\_length）
    \item 段落数（num\_paragraphs）
    \item 平均段落长度（avg\_paragraph\_length）
\end{itemize}

这些特征之所以重要，是因为大语言模型在生成证明时往往具有更机械的段落切分倾向。例如模型更容易每推进一步就另起一段，而人类更倾向于写连续推导。因此“段落数量”最终成为最强判别特征之一。

\paragraph{2. 数学公式与 LaTeX 格式特征}
这一组特征刻画“作者如何书写数学对象”：
\begin{itemize}
    \item 行内公式数量（inline\_math\_count）
    \item 块级公式数量（display\_math\_count）
    \item 数学公式密度（math\_density）
    \item \texttt{\textbackslash frac} 出现次数
    \item 复杂环境数量（latex\_env\_count）
    \item 自适应括号使用（left\_right\_brackets）
    \item 数学字体宏使用（math\_cal\_bb）
    \item 蕴含箭头、QED 结尾符等符号特征
\end{itemize}

这些特征反映了不同来源在公式包装策略上的显著差异。例如，有些模型倾向于把几乎所有变量都写成 \verb|$x$| 形式，有些模型则更频繁使用 \verb|\begin{align}| 这类正规环境，而人类答案往往更“经济”、更随手。

\paragraph{3. 语义与论证风格特征}
这一组特征刻画“作者如何推动论证”：
\begin{itemize}
    \item 逻辑连接词密度（logical\_words\_density）
    \item 祈使句/代词密度（declarative\_density），例如 we, let, suppose, consider
    \item 递进词密度（transition\_words\_density），例如 firstly, moreover, finally
    \item 结论词密度（conclusion\_words\_density）
    \item 枚举结构数量（num\_list\_items）
\end{itemize}

从经验上看，大模型特别喜欢使用“Let us consider...”“We now prove...”“Finally...”这类标准化句式，而人类写作往往更加压缩、跳跃或不那么模板化。这些差异正是我们希望编码进分类器的重要风格信息。

\subsubsection{特征标准化与融合}
自定义特征提取出来后，会先经过 \texttt{StandardScaler} 进行标准化，使不同量纲的特征（例如“段落数”和“逻辑词密度”）能够在统一尺度下进入分类器。随后，标准化后的结构特征与 TF-IDF 稀疏向量进行拼接，共同组成最终输入表示。这意味着最终分类器看到的不是单一模态特征，而是：
\[
\text{final feature} = \text{TF-IDF lexical features} \oplus \text{custom structural features}
\]

\subsection{模型训练}

\subsubsection{为什么要测试多种分类器}
在得到最终特征表示后，我们没有预设“某一种模型一定最好”，而是让多种经典机器学习分类器在相同特征输入下公平比较。原因在于，不同分类器对特征空间的适应能力不同：
\begin{itemize}
    \item 线性模型更适合高维、近似线性可分的文本特征；
    \item 树模型更擅长捕捉非线性边界与特征交互；
    \item 集成模型通常对复杂的手工特征更加敏感。
\end{itemize}

\subsubsection{参与比较的五种模型}
本项目测试了五种代表性算法：
\begin{enumerate}
    \item \textbf{Linear SVM}：作为经典文本分类基线，适合高维稀疏特征；
    \item \textbf{Logistic Regression}：线性概率模型，解释性较好；
    \item \textbf{Random Forest}：袋装树集成，能够处理非线性特征关系；
    \item \textbf{Gradient Boosting}：逐步纠错的提升树模型；
    \item \textbf{Hist Gradient Boosting}：面向更大规模数据和更快训练的高效提升树实现。
\end{enumerate}

\subsubsection{训练与评估协议}
为了分析样本规模对模型性能的影响，我们固定测试集为 200 道题（对应 1000 条记录），并分别设置 20、50、100、400、800 道题作为训练集。每个训练规模下，五个模型都在完全相同的数据划分与特征流水线下训练和测试，最终用测试准确率（accuracy）比较性能。

这个实验设计回答的是一个非常实际的问题：如果真实教学环境中只能获得极少量已标注样本，分类器是否仍然有效？随着训练数据增加，不同模型谁的性能爬升更快？

\section{特征工程分析}
为了更清楚地理解分类器为什么有效，我们并不只关心最终准确率，还关心特征空间中是否真的存在稳定的来源差异。通过随机森林特征重要性分析与 PCA 可视化，我们发现最关键的判别信号主要集中在如下几个方面：
\begin{itemize}
    \item 段落数与平均段落长度：人类与模型在证明组织方式上差异极大；
    \item 行内公式数量：某些模型具有明显的“过度公式包装”倾向；
    \item 祈使句密度：模型在推进证明时更依赖模板化叙述；
    \item 数学公式密度与行数：不同来源在证明展开节奏上表现出稳定差异。
\end{itemize}

\section{实验设置与结果}

\subsection{分类器与数据规模实验}
为了系统比较不同算法在不同数据规模下的性能，我们构建了一个 $5 \times 5$ 的准确率主表：行表示五种分类器，列表示五种训练题量（20、50、100、400、800），测试集始终固定为 200 道题。这样的设计可以同时比较两个维度：
\begin{itemize}
    \item 同一模型随训练数据增加的性能变化；
    \item 同一训练规模下不同模型之间的优劣。
\end{itemize}

\subsection{准确率主表}
表 \ref{tab:results} 给出了完整的实验结果。从结果可以看出，Hist Gradient Boosting 在所有训练规模上都取得了最优或并列最优表现，说明非线性树集成模型最适合处理本项目中“词汇特征 + 结构特征”的混合输入。

\begin{table}[H]
\centering
\caption{不同训练题量下五种分类器的测试准确率（测试集固定为 200 题）}
\label{tab:results}
\begin{tabular}{lccccc}
\toprule
\textbf{模型} & \textbf{20题} & \textbf{50题} & \textbf{100题} & \textbf{400题} & \textbf{800题} \\
\midrule
Hist Gradient Boosting & \textbf{0.868} & \textbf{0.895} & \textbf{0.916} & \textbf{0.947} & \textbf{0.955} \\
XGBoost (sklearn GB) & 0.824 & 0.870 & 0.909 & 0.939 & 0.944 \\
Random Forest & 0.850 & 0.895 & 0.909 & 0.928 & 0.934 \\
Linear SVM & 0.782 & 0.833 & 0.854 & 0.912 & 0.917 \\
Logistic Regression & 0.787 & 0.841 & 0.853 & 0.895 & 0.905 \\
\bottomrule
\end{tabular}
\end{table}

\subsection{可解释性可视化}
除了主表之外，我们还希望说明“为什么这些模型能分得开”。因此报告不再强调准确率折线，而重点展示最佳特征与 PCA 投影图。

图 \ref{fig:pca} 是基于全部结构特征做出的 PCA 二维投影。从中可以观察到，人类解答在特征空间中形成了相对紧密且独立的聚类区域，而不同模型之间则呈现不同程度的重叠。这说明我们构造的特征并不是随意噪声，而是真正提取到了不同来源的稳定风格结构。

图 \ref{fig:violin} 是“段落数”这一最强特征的小提琴图。它直观展示了不同来源在段落切分上的分布差异：模型答案的段落数整体更高、尾部分布更长，而人类答案更集中、更紧凑。

\begin{figure}[H]
    \centering
    \includegraphics[width=0.75\textwidth]{figures/pca_clusters_2d.png}
    \caption{基于结构特征的 PCA 二维投影。可以看到人类解答簇与模型解答簇存在明显分离。}
    \label{fig:pca}
\end{figure}

\begin{figure}[H]
    \centering
    \includegraphics[width=0.75\textwidth]{figures/violin_num_paragraphs.png}
    \caption{“段落数”特征的小提琴图。大语言模型更倾向于频繁分段，而人类解答通常更集中、更连贯。}
    \label{fig:violin}
\end{figure}

\section{结论}
本项目证明了：即使在同一道数学题上，不同来源的解答在内容上高度相似，它们在“如何书写证明”这件事上仍然保留了足够强的来源指纹。通过构建从数据到特征再到分类器的完整流水线，我们发现：
\begin{enumerate}
    \item 数学证明来源识别并不只依赖词汇，而高度依赖结构、公式与论证习惯；
    \item 手工设计的 28 个结构化特征具有很强的可解释性和判别能力；
    \item 树集成模型明显优于线性模型，尤其是 Hist Gradient Boosting 在所有训练规模上都表现最优；
    \item 即便训练样本极少（如 20 题），分类器仍可保持较强识别能力，而在 800 题训练规模下准确率可达到 95.5\%。
\end{enumerate}

因此，这一项目不仅实现了一个高精度的来源分类器，更重要的是提供了一套可解释、可扩展、可迁移的数学文本风格分析框架，可进一步推广到学术作业检测、跨模型写作风格分析以及更广泛的 AI 生成内容取证任务中。

\end{document}
